\chapter*{Abstract}
\addcontentsline{toc}{chapter}{Abstract}

Large Language Models (LLMs) are increasingly used as the reasoning and
code--generating backbone of Automated Machine Learning (AutoML)
pipelines. In such systems an LLM first produces a natural--language
or semi--structured plan that describes how a predictive model will be
built, and then emits executable code that is expected to realise
exactly that plan. Existing evaluations of these pipelines check that
the generated program compiles, runs to completion, and reaches an
acceptable score on some downstream metric. None of these checks
answers a strictly stronger question: does the emitted code
mechanically implement the plan for which the LLM was prepared?

This project formulates the above question as a stand--alone
verification problem that we call \emph{reasoning--code alignment}
verification. The project has three parts, which the report is
careful to keep distinct. First, as a \emph{design} contribution, we
specify a verification framework that lifts an AutoML plan to a set
of atomic typed claims (a Reasoning Intermediate Representation,
RIR), lowers the generated Python program to a directed attributed
multigraph fusing abstract syntax, dataflow and library--level
semantics (a Program Intermediate Representation, PIR), and maps the
two so that each claim receives a three--valued verdict
(\textsc{aligned}, \textsc{misaligned} or \textsc{unverifiable}) and
the pipeline a four--valued one that adds \textsc{partially aligned},
with an accompanying AutoML--specific mismatch taxonomy and a
coverage, weighted--alignment, critical--mismatch and localisation
metric suite.

Second, as an \emph{implementation} contribution, we build a working
end--to--end LLM--driven AutoML stack --- five FastAPI microservices
and two front--ends, a Redis--backed job queue and worker, a
configurable LLM runner (Ollama or OpenRouter), a sandboxed subprocess
executor, and a PostgreSQL / MinIO store --- that generates
scikit--learn pipelines under three prompt
conditions (B0 naive, B1 schema, B2 meta--feature--guided) and runs a
first mechanical faithfulness verifier on every B2 run, checking that
each cited action appears in the generated code and that each cited
dataset statistic matches the real meta--feature.

Third, as an \emph{empirical} contribution, we evaluate the stack on
a controlled 270--cell sweep (6 tabular datasets $\times$ 3
conditions $\times$ 3 model backends $\times$ 5 seeds). The naive and
schema conditions are completely reliable (100\,\% runnable code),
whereas the meta--feature condition is markedly less reliable
(62.2\,\% success) and 5--25$\times$ more expensive in prompt tokens,
without any statistically significant improvement in predictive score
(paired Wilcoxon, all $p>0.09$). Crucially, the faithfulness verifier
flagged seven B2 runs that executed successfully and produced good
scores but whose stated reasoning did not match their code --- a
failure mode that execution-- and performance--based checks cannot
detect. The richer graph--based verifier, the independent pre--code
reasoning capture and the baseline comparison are designed --- with a
rule--based verifier prototype on a feature branch --- but not yet
evaluated at population scale; the report marks these as future work
rather than completed results.

\medskip

\noindent\textbf{Keywords:} automated machine learning, large language
models, LLM agents, code generation, program analysis, mechanical
verification, reasoning--code alignment, reasoning faithfulness.
